Fix \(P\in\mathfrak P_{\mathcal D}^{\mathrm{obs}}\).  For the arm-normalized profile, positivity and finiteness follow from positive unit costs, Assumption~\(\mathrm{ass{:}nondegenerate\text{-}components}\), and Theorem~\(\mathrm{thm{:}budget\text{-}profile}\).  For the IKM-normalized profile, Assumption~\(\mathrm{ass{:}source\text{-}arm\text{-}positivity}\) and Definition~\(\mathrm{def{:}ikmw\text{-}normalized\text{-}design\text{-}value}\) give
\[
 0<(1-p_{g,1})^2V_{g,d}<\infty
 \qquad(g\in\{E,O\},\ d\in\mathcal D),
\]
so positive costs again imply \(0<w_d<\infty\) for every \(d\).  Thus the following argument applies to either displayed profile.

Define the finite treatment-aware order relation
\[
 \mathcal O_P^A
 =\{(d,d')\in\mathcal D^2:R_d^A(P)\le R_{d'}^A(P)\}.
\]
It is nonempty because it contains every diagonal pair.  If \(C\) satisfies all treatment-aware order implications, then division by \(w_{d'}>0\) gives
\[
 \frac{w_d}{w_{d'}}\le C
 \qquad\text{for every }(d,d')\in\mathcal O_P^A.
\]
Consequently \(C\) is at least the maximum in (1).  Conversely, that finite maximum dominates every displayed ratio, so it satisfies every implication.  This proves (1) and its sharpness.

Choose a design minimizing \(w_d\) among \(\mathcal D_{\mathrm{pred}}^A(P)\) and a design minimizing \(w_d\) over all of \(\mathcal D\).  The first design minimizes \(R_d^A(P)\), so the pair consisting of these two choices belongs to \(\mathcal O_P^A\).  Hence the ratio of their profile values is at most (1).  By Definitions~\(\mathrm{def{:}predictive\text{-}benchmark}\) and \(\mathrm{def{:}ikmw\text{-}normalized\text{-}design\text{-}value}\), that ratio is respectively \(\rho_{\mathrm{pred}}^A(P)\) or \(\rho_{\mathrm{pred}}^{A,\mathrm{IKM}}(P)\).  Division by the strictly positive global minimum also shows that each ratio is the least constant for its actual best-tie selector and belongs to \([1,\infty)\).  Taking suprema over a nonempty law class proves that a finite uniform best-tie constant exists if and only if the supremum of the corresponding pointwise ratio is finite.  This is the treatment-aware analogue of the finite-order argument in Theorem~\(\mathrm{thm{:}sharp\text{-}predictive\text{-}operator\text{-}certificate}\).

Now impose the stated treatment-aware quadratic-form realization, with all infima and suprema taken over the cones supplied by \(\mathsf H_{\mathrm{comp}}\).  The envelope definitions give, for every admissible \(d\) and direction \(x\),
\[
 \Lambda_-^A(P)Q_{R,d}^A(x)
 \le F_d(x)
 \le \Lambda_+^A(P)Q_{R,d}^A(x).                           \tag{10}
\]
For every \((d,d')\in\mathcal O_P^A\), evaluation at the transported law-specific directions yields
\[
\begin{aligned}
 \mathcal V_d^*(P)
 &=F_d\{x_d(P)\}\\
 &\le\Lambda_+^A(P)R_d^A(P)
 \le\Lambda_+^A(P)R_{d'}^A(P)\\
 &\le\frac{\Lambda_+^A(P)}{\Lambda_-^A(P)}
 F_{d'}\{x_{d'}(P)\}\\
 &=\frac{\Lambda_+^A(P)}{\Lambda_-^A(P)}
 \mathcal V_{d'}^*(P).
\end{aligned}                                               \tag{11}
\]
Maximizing (11) over \(\mathcal O_P^A\), and using the best-tie bound just proved, gives (2).

Finally, Theorem~\(\mathrm{thm{:}sharp\text{-}predictive\text{-}operator\text{-}certificate}\) proves the profile identity
\[
 F_d(x)=\inf_{0<\alpha<1}
 \left\{
 \frac{c_{E,d}Q_{E,d}(x)}{\alpha}
 +\frac{c_{O,d}Q_{O,d}(x)}{1-\alpha}
 \right\}
\]
and gives an explicit two-dimensional example in which \(F_d\) violates the parallelogram identity.  Thus the profiled numerator is not generally quadratic.  A largest generalized eigenvalue at a fixed allocation supplies at most a one-sided full-space upper envelope; it does not provide the positive lower envelope in (10).  Even when both envelopes exist, equality in (11) additionally requires law-realized treatment-aware-risk-ordered directions to attain the relevant extrema.  Therefore the envelope is generally conservative and the sharp treatment-aware selection constant is not generally a largest generalized eigenvalue.